% -------------------------------------------------------------------------------------------------
%
%  Skeleton for semester and master thesis reports at the Laboratory for Software Technology.
%  Based on the skeleton provided by the Institute of Robotics and Intelligent Systems.
%
% -------------------------------------------------------------------------------------------------
%
% USAGE:      compile with PDFLaTeX
%
% HISTORY:    - written by Sascha A. Stoeter <stoeter@iris.ethz.ch>, www.stoeter.com, 02.06.2004
%             - modified by Martin Probst, 18.08.2004
%             - extended and adapted for use at LST by Oliver Trachsel, 2007-08-21
% -------------------------------------------------------------------------------------------------
\documentclass[11pt,a4paper]{book}
\usepackage{lstreport}

% -------------------------------------------------------------------------------------------------
% Add needed packages. Some generally useful packages are listed for
% your convenience.
% -------------------------------------------------------------------------------------------------
\usepackage{subfigure}                          % enable the use of subfigures
\usepackage[thickspace,thinqspace]{SIunits}     %
\usepackage[plainpages=false,pdfpagelabels]{hyperref}    % enable hyperlinks in pdf/ps Docs
\usepackage{listings}                           % to embed source code

% -------------------------------------------------------------------------------------------------
% Select type of thesis
% -------------------------------------------------------------------------------------------------
\renewcommand{\thesistype}{Semester}
%\renewcommand{\thesistype}{Diploma}
%\renewcommand{\thesistype}{Master}

% -------------------------------------------------------------------------------------------------
% Set names
% -------------------------------------------------------------------------------------------------
\renewcommand{\thesisauthor}{Peter Muster}
\renewcommand{\thesisadvisor}{Hugh Watts}

% -------------------------------------------------------------------------------------------------
% Beginning of the main document body
% -------------------------------------------------------------------------------------------------
\begin{document}

% -------------------------------------------------------------------------------------------------
% Front matter with title page, table of contents, and abstracts 
% -------------------------------------------------------------------------------------------------
\frontmatter

% Title page: set title and date 
\thesistitlepage{Thesis Title}{August 2007}

% Abstract must not be longer than one page per language. English and
% German abstracts are mandatory.
\chapter*{Abstract}
Short summary of thesis.

\chapter*{Zusammenfassung}
Kurzfassung der Arbeit.

% Table of contents
\tableofcontents

% -------------------------------------------------------------------------------------------------
% Main document body
% -------------------------------------------------------------------------------------------------
\mainmatter
\chapter{Introduction}
\label{s:Introduction}
Explain scope and structure of report.

\chapter{Great Work}
\label{s:GreatWork}

This and the following chapters detail the original work.

\chapter{Examples}
\label{s:Examples}

This chapter provides some additional hints and examples for the
layout and style of the thesis. It is worthwhile to look at the source
file \verb|Examples.tex| for this appendix to understand how it was
created.

\section{Tables}

Tables are left justified and the caption appears on top as seen in
Table~\ref{t:Translations}.

\begin{table}[ht]
\centering
\begin{tabular}{ll}
\hline
\textbf{English} & \textbf{German}\\
\hline
cell phone       & Handy\\
Diet Coke        & Coca Cola light\\
\hline
\end{tabular}
\caption[Translations]{\label{t:Translations}Translations.}
\end{table}

\section{Figures}

Figure~\ref{f:IRISlogo} shows a simple figure with a single picture
and Figure~\ref{f:SubfigureExample} shows a more complex figure
containing subfigures.

\begin{figure}[ht]
\centering
\includegraphics[width=.6\linewidth]{figures/LSTlogo}
\caption[LST logo]{\label{f:IRISlogo}LST logo.}
\end{figure}

\begin{figure}[ht]
\centering
\subfigure[ETH logo]{\includegraphics[height=12mm]{figures/ETHlogo}}\quad
\subfigure[LST logo]{\includegraphics[height=12mm]{figures/LSTlogo}}
\caption[Subfigure example]{\label{f:SubfigureExample}Two pictures as
  part of a single figure through the magic of the subfigure package.}
\end{figure}

\section{Units}

The SIUnits package provides nice spacing for units as demonstrated in
Table~\ref{t:SIUnits}. Use of the package also makes it easy to change
the style or even the unit text in the future.

\begin{table}[ht]
\centering
\begin{tabular}{ll}
\hline
\textbf{Output}   & \textbf{Command}\\
\hline
42m               & \verb|42m|\\
\unit{42}{\metre} & \verb|\unit{42}{\metre}|\\
42 m              & \verb|42 m|\\
\hline
\end{tabular}
\caption[Spacing for units]{\label{t:SIUnits}Spacing for units.}
\end{table}

\section{Source code}

The listings package provides tools to typeset source code
listings. It supports many programming languages and provides a lot of
formatting options.

\lstset{numbers=left, numberstyle=\tiny, stepnumber=1, numbersep=5pt}
\lstset{basicstyle=\ttfamily}
\lstset{frame=tb}

\begin{lstlisting}[float,caption=Example usage of the listing package,label=l:javaClass,language=Java]
class S {
   int f1 = 42;
   public S(int x) {
          f1 = x;
   }
}
\end{lstlisting}

Listing \ref{l:javaClass} shows an example listing. Code snippets can
also be inserted in normal text:
\verb$\lstinline|int f1 = 42;|$ gives \lstinline$int f1 = 42;$

\section{Miscellany}

\begin{description}

\item[Capitalization.] When referring to a named table (such as in the
  previous section), the word \emph{table} is capitalized. The same is
  true for figures, chapters and sections.

\item[Bibliography.] Use \verb|bibtex| to make your life easier and to
  produce consistently formatted entries.

\item[Contractions.] Avoid contractions. For instance, use ``do not''
  rather than ``don't.''

\item[Style guide.] A classic reference book on writing style is
  Strunk's \emph{The Elements of Style} \cite{Strunk-ElementsOfStyle}.

\end{description}

% -------------------------------------------------------------------------------------------------
% Appendices (if needed)
% -------------------------------------------------------------------------------------------------
\appendix
\chapter{Extra Stuff}
\label{s:ExtraStuff}

Additional material such as long mathematical derivations.

% -------------------------------------------------------------------------------------------------
% Bibliography
% -------------------------------------------------------------------------------------------------
\addcontentsline{toc}{chapter}{Bibliography}
\bibliography{bibliography}

\end{document}
